%% ============================================================
%% 分片 07：原文第 3 节「Topological Field Theory From ${\EUN}=4$ Super Yang-Mills Theory」（原文 glshort.tex 起于第 1216 行）
%% 由 tools/harvmac2tex.py 从 arXiv 源码（hep-th/0604151v3）自动生成。
%% 只译英文散文；公式、编号、\label、\cite 一律不改。
%% ============================================================

\subsection{拓扑拉格朗日量}

\label{sec:toplag}

现在我们的目标是找到一个拉格朗日量：它对任意 $t$ 值都具有拓扑对称，并且在 $M$ 平坦时退化为底层 ${\EUN}=4$ 超对称杨--米尔斯理论的拉格朗日量。我们可以手工来做这件事——从平坦 $M$ 下的式~\eqref{eq:ygo} 出发，追问当 $M$ 不再平坦时，需要哪些依赖于曲率的项才能保持拓扑对称。但我们将采取另一种途径。

一个非常有用的第一步，是计算超对称 $Q_\ell$ 与 $Q_\rr$ 在场上生成的代数；等价地，计算 $Q=uQ_\ell+vQ_\rr$ 生成的代数。对场 $A,\phi,\psi$、$\tilde\psi$ 与 $\sigma$，我们算得
% 式 (3.35)
\setcounter{equation}{34}
\begin{equation}\label{eq:ibo}
\begin{aligned}\delta_T{}^2 A & = -i(u^2+v^2)(-D\sigma)\\
\delta_T{}^2\phi & = -i(u^2+v^2)[\sigma,\phi]\\ \delta_T{}^2 \psi
& = -i(u^2+v^2)[\sigma,\psi]\\ \delta_T{}^2\tilde\psi& =
-i(u^2+v^2)[\sigma,\tilde\psi]\\ \delta_T{}^2\sigma& = 0.\\\end{aligned}
\end{equation}

这些结果可以概括为：若 $\Phi$ 是任何 $\EUK\geq 0$ 的场，则 $\delta_T{}^2\Phi=-i(u^2+v^2){{\lie}}_\sigma(\Phi)$，其中 ${{\lie}}_\sigma(\Phi)$ 是 $\Phi$ 在由 $\sigma$ 生成的规范变换下的一阶变化（因此 $\lie$ 代表 Lie 导数）。例如 ${{\lie}}_\sigma(A)=-D\sigma$；对此处的其他场，${{\lie}}_\sigma(\Phi)=[\sigma,\Phi]$。关系式 $\delta_T{}^2\Phi=-i(u^2+v^2){{\lie}}_\sigma(\Phi)$ 可以按 $u$ 与 $v$ 的幂展开，它等价于
% 式 (3.36)
\setcounter{equation}{35}
\begin{equation}\label{eq:uvu}
\begin{aligned}\delta_\ell^2\Phi=\delta_r^2\Phi&=-i{{\lie}}_\sigma(\Phi)\\
\{\delta_\ell,\delta_\rr\}\Phi& = 0.\\\end{aligned}
\end{equation}

然而，若对 $\EUK<0$ 的场计算 $\delta_T{}^2\Phi$，并不能以显然的方式得到 $-i(u^2+v^2){{\lie}}_\sigma(\Phi)$。事实上，在扭化后的 ${\EUN}=4$ 超对称杨--米尔斯理论中，关系式 $\delta_T{}^2\Phi=-i(u^2+v^2){{\lie}}_\sigma(\Phi)$ 对所有 $\Phi$ 都成立；但对某些场，它只有在用到运动方程时才成立。

\begin{yizhu}[易错点]
按：(3.35) 给出的代数对 $\EUK\geq 0$ 的场无需运动方程即成立。$\EUK$ 是扭化所用的那个 $U(1)$ 荷（第~\ref{sec:twisting}~节已逐个标定各场取值）：$A$ 与 $\phi$ 的 $\EUK=0$，$\psi$ 与 $\tilde\psi$ 的 $\EUK=1$，$\chi,\eta,\tilde\eta$ 的 $\EUK=-1$，而 $\sigma$ 的 $\EUK=2$、$\bar\sigma$ 的 $\EUK=-2$；相应地，$\delta_T$ 每作用一次便把 $\EUK$ 提高 1。$\sigma$ 位于最高处，$[\sigma,\sigma]=0$ 使 (3.35) 的最后一行 $\delta_T{}^2\sigma=0$ 与通式自洽；反过来，对 $\EUK<0$ 的 $\chi,\eta,\tilde\eta$，其变换含 $\sigma$ 以外的非线性项，代数在壳外并不闭合，这正是下一段需要辅助场的根源。
\end{yizhu}

一般说来，构造超对称拉格朗日量要容易得多，只要能够引入「辅助场」——它们是额外的场，若愿意，最终可以用其运动方程消去；引入它们的目的是让代数（在这里即 $\delta_T{}^2\Phi=-i(u^2+v^2){{\lie}}_\sigma(\Phi)$）无需借助运动方程便得到满足。在具有大量超对称的理论（例如 ${\EUN}=4$ 超对称杨--米尔斯理论）中，要找到合适的辅助场可能非常困难。正如 \cite{ref53} 首先指出的那样，本例正是这一问题出现的情形。

不过，这一困难可以通过引入一个辅助场 $P$ 而大大缓解：$P$ 是取值于李代数的零形式，引入它是为了让 $\eta$ 与 $\tilde\eta$ 的代数闭合（但不包括 $\chi$）。我们取 $\bar\sigma,\eta,\tilde\eta$ 与 $P$ 的变换律为
% 式 (3.37)
\setcounter{equation}{36}
\begin{equation}\label{eq:yfog}
\begin{aligned}\delta_T\bar\sigma & = iu\eta +i v\tilde\eta
\\
                    \delta_T\eta & = vP+u[\bar\sigma,\sigma]\\
                    \delta_T\tilde\eta & =
                    -uP+v[\bar\sigma,\sigma]\\
                    \delta_T P & =
                    -iv[\sigma,\eta]+i u[\sigma,\tilde\eta].\\\end{aligned}
\end{equation}

这些方程一旦施加运动方程，就化为第~\ref{sec:twisting}~节中的那些方程——值得注意的是，它们会令 $P=D^*\phi$。

一旦在某一组场上得到闭合代数——在本例中，即理论中除 $\chi$ 之外的所有场（含 $P$），因为我们尚未让代数在 $\chi$ 上闭合——寻找可能的不变拉格朗日量一般是件简单的事。在本例中，我们希望为代数已在其上闭合的那些场找到一个部分作用量 ${{\cmmib I}}_0$，它满足 $\delta_\ell{{\cmmib I}}_0=\delta_\rr{{\cmmib I}}_0=0$，从而对任意 $u,v$ 满足 $\delta_T{{\cmmib I}}_0=0$，其中 $\delta_T=u\delta_\ell+v\delta_\rr$。只要 ${{\cmmib I}}_0=\delta_\ell\delta_\rr V$（$V$ 为某个规范不变者），这些性质便成立。确实，若 $V$ 是规范不变的，则 ${{\lie}}_\sigma(V)=0$，而式~\eqref{eq:uvu} 便化为 $\delta_\ell^2V=\{\delta_\ell,\delta_\rr\}V=\delta_\rr^2V=0$。于是 $\delta_\ell\delta_\rr V$ 自动被 $\delta$ 湮灭。

此外，我们希望选取 $V$，使得 $P$ 的运动方程为 $P=D^*\phi$，以便与第~\ref{sec:twisting}~节中的公式相容。因此我们取% 式 (3.38)
\setcounter{equation}{37}
\begin{equation}\label{eq:obo}
V={2\over e^2}\int_M
d^4x\sqrt g\left(-{1\over 2}\Tr\,\eta\tilde\eta -i\Tr\,\bar\sigma
D^*\phi\right),
\end{equation}
并算得 ${\cmmib I}_0=\delta_\ell\delta_\rr V$ 为% 式 (3.39)
\setcounter{equation}{38}
\begin{equation}\label{eq:nobo}
\begin{aligned}{{\cmmib
I}}_0= {2\over e^2}\int_Md^4x\sqrt g\Tr& \left({1\over 2}P^2
-PD^*\phi+{1\over 2}[\bar\sigma,\sigma]^2-
D_\mu\bar\sigma\,D^\mu\sigma
-[\phi_\mu,\sigma][\phi^\mu,\bar\sigma]\right.\\ & + i\tilde\eta
D_\mu\tilde\psi^\mu+i\eta
D_\mu\psi^\mu-i\tilde\eta[\psi_\mu,\phi^\mu]+i\eta[\tilde\psi_\mu,\phi^\mu]
\\ &\left.-{i\over 2}[\sigma,\tilde\eta]\tilde\eta -{i\over
2}[\sigma,\eta]\eta+i[\bar\sigma,\psi_\mu]\psi^\mu
+i[\bar\sigma,\tilde\psi_\mu]\tilde\psi^\mu\right).\\\end{aligned}
\end{equation}
$P$ 的 Euler--Lagrange 方程正如所愿为 $P=D^*\phi$，并且可以作替换% 式 (3.40)
\setcounter{equation}{39}
\begin{equation}\label{eq:olgog}
\Tr\left(
P^2-2PD^*\phi\right)\to-\Tr (D^*\phi)^2,
\end{equation}
若愿意的话。

${{\cmmib I}}_0$ 显然对任意 $t=v/u$ 都具有拓扑对称。此外，由于它具有 $\delta_\ell\delta_\rr V$ 的形式——其中 $M$ 的度规并不进入 $\delta_\ell$ 与 $\delta_\rr$ 的定义，而只进入 $V$ 的选取——它还具有导致拓扑场论的关键性质：它对 $M$ 的度规的依赖具有 $\delta_\ell\delta_\rr(\dots)$ 的形式。

\begin{yizhu}[物理图像]
按：这就是拓扑场论的定义方式之一：作用量对度规的依赖必须是 $\delta_t$ 恰当的形式。由于 $\delta_t$ 与度规无关，把度规从 $g$ 形变到 $g'$ 时，作用量的改变可以写成 $\delta_t(\dots)$，于是对 $\delta_t$-闭的可观测量，其关联函数在度规形变下不变——$\delta_t$-恰当项在 $Q$-上同调意义下为零。要注意这里说的是 $\delta_tV$ 中「看得见」的度规依赖；$\delta_t$ 本身的定义不含度规，这一点在把理论放到弯曲 $M$ 上时至关重要。参见 Witten 2008 综述中关于拓扑扭化的一般讨论。
\end{yizhu}

${{\cmmib I}}_0$ 之所以不是令人满意的作用量，唯一原因在于它退化：它既不含规范场的动能，也不含任何含 $\chi$ 的项。要想以像上面那样简单的构造来完成作用量的构造，似乎做不到。虽然可以添加辅助场使代数在 $\chi$ 上闭合，但似乎无法以同时纳入 $\delta_\ell$ 与 $\delta_\rr$、并且有助于理解相应的扭化 ${\EUN}=4$ 作用量的方式来做到这一点。（\eqref{eq:ormob} 的右端出现一个依赖于 $t$ 的拓扑不变量，看来正是此事的障碍。）于是我们固定一个特定的 $t$ 值，只考虑微分 $\delta_t=\delta_\ell+t\delta_\rr$。我们添加一个辅助场 $H$，它将是取值于李代数的 2-形式，并假设% 式 (3.41)
\setcounter{equation}{40}
\begin{equation}\label{eq:hco}
\begin{aligned}\delta_t\chi& = H\\
                           \delta_t H & =
                           -i(1+t^2)[\sigma,\chi].\end{aligned}
\end{equation}

这与
% 式 (3.42)
\setcounter{equation}{41}
\begin{equation}\label{eq:kino}
\delta_t^2(\Phi)=-i(1+t^2){{\lie}}_\sigma(\Phi),
\end{equation}
相容（$\Phi=\chi,H$）；它是 $\delta_T^2=-i(u^2+v^2)\lie_\sigma$ 在 $(u,v)=(1,t)$ 处的特化。为使式~\eqref{eq:hco} 与第~\ref{sec:twisting}~节中得到的 $\chi$ 的变换相一致，运动方程必须给出 $H^+=\EUV^+(t)$、$H^-=t\EUV^-(t)$。我们将构造作用量来保证这一点。

与之前一样，被 $\delta_t$ 湮灭的作用量可以取 ${\cmmib I}_1=\delta_t V_1$，其中 $V_1$ 为任意规范不变者。我们取
% 式 (3.43)
\setcounter{equation}{42}
\begin{equation}\label{eq:usefuld}
V_1={2\over e^2}\int_Md^4x \sqrt g{1\over
(1+t^2)}\Tr\left(\chi^+_{\mu\nu}\left({1\over
2}H^{+\mu\nu}-\EUV^+(t)^{\mu\nu}\right)
+\chi^-_{\mu\nu}\left({1\over
2}H^{-\mu\nu}-t\EUV^-(t)^{\mu\nu}\right)\right).
\end{equation}
于是
% 式 (3.44)
\setcounter{equation}{43}
\begin{equation}\label{eq:usefuldud}
\begin{aligned}{{{\cmmib I}}}_1= {1\over e^2}\int_Md^4x
\sqrt g&\left({2\over (1+t^2)}\Tr\left({1\over
2}H^+_{\mu\nu}H^{+\mu\nu}-H^+_{\mu\nu}\EUV^+(t)^{\mu\nu}\right)\right.\\
& + {2\over (1+t^2)}\Tr\left({1\over
2}H^-_{\mu\nu}H^{-\mu\nu}-tH^-_{\mu\nu}\EUV^-(t)^{\mu\nu}\right)\\
&+
2\Tr\left(\chi^+_{\mu\nu}\left(iD\psi+i[\tilde\psi,\phi]\right)^{\mu\nu}
+\chi^-_{\mu\nu}\left(iD\tilde\psi-i[\psi,\phi]\right)^{\mu\nu}\right)\\
&+\Tr\left(i\chi^+_{\mu\nu}[\sigma,\chi^{+\mu\nu}]+
i\chi^-_{\mu\nu}[\sigma,\chi^{-\mu\nu}] \right)\biggr).\\\end{aligned}
\end{equation}
把 $H^{\pm}$ 从其中积掉，便可写出等价的作用量% 式 (3.45)
\setcounter{equation}{44}
\begin{equation}\label{eq:usefuldudd}
\begin{aligned}{{\cmmib I}}_1= {1\over
e^2}\int_Md^4x \sqrt g&\left(-{t^{-1}\over
(t+t^{-1})}\Tr\,\EUV^+(t)_{\mu\nu}\EUV^+(t)^{\mu\nu}-{t\over
(t+t^{-1})}\Tr\,\EUV^-(t)_{\mu\nu}\EUV^-(t)^{\mu\nu}\right)\\ &+
2\Tr\left(\chi^+_{\mu\nu}\left(iD\psi+i[\tilde\psi,\phi]\right)^{\mu\nu}
+\chi^-_{\mu\nu}\left(iD\tilde\psi-i[\psi,\phi]\right)^{\mu\nu}\right)\\
&+\Tr\bigl(i\chi^+_{\mu\nu}[\sigma,\chi^{+\mu\nu}]+
i\chi^-_{\mu\nu}[\sigma,\chi^{-\mu\nu}] \bigr)\biggr).\\\end{aligned}
\end{equation}

\begin{yizhu}[易错点]
按：$H$ 是 2-形式（其自对偶部分 $H^+$ 与反自对偶部分 $H^-$ 分别与 $\chi^+$、$\chi^-$ 配对），而前面那个 $P$ 是零形式；两者都是为使代数在壳外闭合而引入的辅助场，最后分别由代数运动方程消去。特别留意 (3.43) 中 $H^\pm$ 前的系数并不对称：$H^+$ 配 $\EUV^+(t)$，而 $H^-$ 配 $t\EUV^-(t)$。这个不对称正是 $\delta_t=\delta_\ell+t\delta_\rr$ 中 $t$ 的权重造成的（$\chi^+$ 与 $\chi^-$ 在 $\delta_\ell,\delta_\rr$ 下的相对权重不同）；若想当然地把两项都写成同一个系数，拓扑对称就会被破坏。
\end{yizhu}

最后，有用的恒等式 \eqref{eq:ormob} 使我们能够写出
% 式 (3.46)
\setcounter{equation}{45}
\begin{equation}\label{eq:nuseful}
\begin{aligned}{{\cmmib I}}_1=-{1\over e^2}\int_Md^4x \sqrt
g&\left(\Tr\left({1\over 2}F_{\mu\nu} F^{\mu\nu}+D_\mu\phi_\nu
D^\mu\phi^\nu +R_{\mu\nu}\phi^\mu\phi^\nu+ {1\over
2}[\phi_\mu,\phi_\nu]^2-(D^*\phi)^2
  \right)\right.\\ &-
2\Tr\left(\chi^+_{\mu\nu}\left(iD\psi+i[\tilde\psi,\phi]\right)^{\mu\nu}
+\chi^-_{\mu\nu}\left(iD\tilde\psi-i[\psi,\phi]\right)^{\mu\nu}\right)\\
&-\Tr\left(i\chi^+_{\mu\nu}[\sigma,\chi^{+\mu\nu}]+
i\chi^-_{\mu\nu}[\sigma,\chi^{-\mu\nu}] \right)\biggr)\\ &
~~~~~~~+{t-t^{-1}\over e^2(t+t^{-1})}\int_M\Tr\, F\wedge F.\\\end{aligned}
\end{equation}

除拓扑项之外，${{\cmmib I}}_0+{{\cmmib I}}_1$ 中仅依赖于 $A$ 与 $\phi$ 的那部分也可以写成 \cite{ref53}% 式 (3.47)
\setcounter{equation}{46}
\begin{equation}\label{eq:wingo}
{{\cmmib I}}^{(A,\phi)}=-{1\over
e^2}\int d^4x\,\sqrt g\, \Tr\left({1\over 2}
\cF_{\mu\nu}\bar\cF^{\mu\nu}+(D^*\phi)^2\right),
\end{equation}
正如用一些与证明消灭定理时所用的类似的分部积分可以说明的那样。含 $\sigma$ 的类似项可在式~\eqref{eq:nobo} 中找到，它们为% 式 (3.48)
\setcounter{equation}{47}
\begin{equation}\label{eq:uwingo}
{\cmmib
I}^\sigma={2\over e^2}\int d^4x \sqrt g\,\Tr\left( {1\over
2}[\bar\sigma,\sigma]^2- D_\mu\bar\sigma\,D^\mu\sigma
-[\phi_\mu,\sigma][\phi^\mu,\bar\sigma]\right).
\end{equation}

\begin{yizhu}[含糊点拨]
按：(3.47) 中的 $\cF_{\mu\nu}$ 是原文只用这一次的记号，应理解为把 (3.46) 中只含 $A,\phi$ 的玻色项重新打包之后的复合曲率，而不是 $F_{\mu\nu}$ 的笔误。原文紧接着说，这个等式要「用一些与证明消灭定理时类似的分部积分」才能证明，因此逐项对照 (3.46) 与 (3.47) 是行不通的：$R_{\mu\nu}\phi^\mu\phi^\nu$ 这类项正是在分部积分中被重新分配进 $\cF\bar\cF$ 与 $(D^*\phi)^2$ 的组合里。读 (3.47) 时只需把它当作一个方便的整体写法，其含义由右端给出。
\end{yizhu}

关键的一点——我们在证明消灭定理时已经用过它——是：除最后一项之外，${{\cmmib I}}_1$ 与 $t$ 无关。但最后那一项是拓扑不变量，它自动被 $\delta_t$ 湮灭，而且事实上单独就被 $\delta_\ell$ 与 $\delta_\rr$ 两者湮灭。因此，在不破坏拓扑对称的前提下，我们可以给作用量再加一项% 式 (3.49)
\setcounter{equation}{48}
\begin{equation}\label{eq:ucno}
{{\cmmib I}}_2=-\left({t-t^{-1}\over
e^2(t+t^{-1})}-i{\theta\over 8\pi^2}\right)\int_M\Tr\,F\wedge F.
\end{equation}

我们只是选取系数以抵消 ${{\cmmib I}}_1$ 中的 $t$ 依赖，结果这一项在作用量中留下的是标准的 $\theta$ 参数。

于是，总作用量为 ${{\cmmib I}}={{\cmmib I}}_0+{{\cmmib I}}_1+{{\cmmib I}}_2$。这一构造使下一点变得显然：对构造 ${{\cmmib I}}_1$ 时所用的那个特定 $t$ 值，${{\cmmib I}}$ 被 $\delta_t$ 湮灭。但由于 ${{\cmmib I}}$ 实际上与 $t$ 无关，它对所有 $t$ 都被 $\delta_t$ 湮灭。

此外，${{\cmmib I}}$ 是拓扑场论的作用量。当 $M$ 的度规改变时，${{\cmmib I}}$ 的变化具有 $\delta_t(\dots)$ 的形式：因为 ${{\cmmib I}}_0$ 与 ${{\cmmib I}}_1$ 甚至在变更度规之前就已具有这种形式，而 ${{\cmmib I}}_2$ 完全不含 $M$ 的度规。

我们已完成目标：构造出一个具有完整 ${\Bbb{ CP}}^1$ 族拓扑对称的作用量——这尤其清楚地表明，该理论在闭四维流形 $M$ 上（无算子插入时）的配分函数与 $t$ 无关。容易验证：在通过式~\eqref{eq:olgog} 消去 $P$（此时 $(D^*\phi)^2$ 各项相消）之后，${{\cmmib I}}$ 的玻色部分在 $M=\Bbb{R}^4$ 时化为式~\eqref{eq:loopy}。拓扑对称蕴含费米子项也一致。因此，这个理论是 ${\EUN}=4$ 超对称杨--米尔斯理论到弯曲四维流形的推广。

\subsection{典范参数}

\label{sec:canonpar}

迄今为止，我们得到了一族依赖于耦合参数 $\tau=\theta/2\pi+4\pi i/e^2$ 与扭化参数 $t$ 的拓扑场论。但这些参数并非真的全都独立。下面我们将说明，该理论其实只依赖于 $t$ 与 $\tau$ 的某个特定组合 $\Psi$，我们称之为典范参数。

我们把 ${{\cmmib I}}_2$ 中拓扑项的系数——其形式为 $-\left({(t-t^{-1})/e^2(t+t^{-1})}-i{\theta/8\pi^2}\right)$——写成 $-\Psi/4\pi i$，其中
% 式 (3.50)
\setcounter{equation}{49}
\begin{equation}\label{eq:yorgo}
\Psi = {\tau+\bar\tau\over 2}+{\tau-\bar\tau\over
2}\left({t-t^{-1}\over t+t^{-1}}\right).
\end{equation}
我们称 $\Psi$ 为典范参数。

\begin{yizhu}[含糊点拨]
按：$\Psi$ 与 $\tau$ 不是一回事，前者把 $t$ 混了进来。把 (3.50) 中的 $\tau$ 写成 $\tau={\rm Re}\,\tau+i\,{\rm Im}\,\tau$（注意 ${\rm Im}\,\tau=4\pi/e^2>0$），便有
\[
\Psi={\rm Re}\,\tau+i\,{\rm Im}\,\tau\cdot{t-t^{-1}\over t+t^{-1}}.
\]
由此立刻看出本节后面列举的几个特殊值：$|t|=1$ 时 $(t-t^{-1})/(t+t^{-1})$ 为纯虚数，于是 $\Psi$ 为实；$|t|>1$ 时 $\Psi$ 落在上半平面（$|t|<1$ 时落在下半平面）；$t=\pm i$ 时分母为零，$\Psi=\infty$，与 $\tau$ 无关。也就是说 $\Psi$ 的取值范围是整个 $\Bbb{CP}^1$，$\tau$ 只控制 $\Psi$ 的实部与虚部的尺度。
\end{yizhu}

我们断言，任何可观测量 ${\EUO}(\tau,\bar\tau,t)$ 其实只依赖于典范参数 $\Psi$。确实，该理论的作用量具有形式% 式 (3.51)
\setcounter{equation}{50}
\begin{equation}\label{eq:juty}
{{\cmmib I}}=\delta_tV
+{i\Psi\over 4\pi }\int_M\Tr \,F\wedge F.
\end{equation}
在 $\delta_tV$ 项无关紧要的程度上，理论自然只依赖于系数 $\Psi$。

乍看之下，这一论证有个谬误。除了 ${{\cmmib I}}$ 中可能显含的 $t$ 依赖（我们已经安排得没有这种依赖，尽管各个项分别依赖于 $t$）之外，还可能存在一种 $t$ 依赖，它来自 $\delta_t$ 的定义本身依赖于 $t$。毕竟，理论的界定既依赖于作用量，也依赖于拓扑对称 $\delta_t$。但事实上，只要 $t\not=\pm i$，就可以通过重新定义场来消除 $\delta_t$ 的 $t$ 依赖。把 $\delta_t$ 重标度为 $\delta'_t=\delta_t/\sqrt{1+t^2}$，于是代数变为 $(\delta'_t)^2=-i{{\lie}}_\sigma$，与 $t$ 无关。然后，在把辅助场 $P$ 与 $H$ 包含进来、并对费米子作一些小的重新定义之后，每个多重态都取形式 $\delta'_t U=V$、$\delta'_tV=-i{{\lie}}_\sigma(U)$，其中 $U,V$ 为某对场。（这是上同调场论中标准的一类多重态，在数学上与等变上同调的 Cartan 模型相关。）这样一来，多重态的变换律就完全与 $t$ 无关了。

\par\medskip{\it 若干特殊值}

关于 $\Psi$、$t$ 与 $\tau$ 之间关系的一些特殊情况值得注意。我们尤其关注那些在 $S$-对偶下行为简单、并且在紧化到二维之后解释简单的取值。

若 $t=\pm 1$，则 $\Psi={\rm Re}\,\tau$。因此，任何实值的 $\Psi$ 都可以在 $t=\pm 1$ 处达到，只需改变规范场论的 $\theta$ 角；但这时只可能出现实值。更一般地，若 $|t|=1$，则 $\Psi$ 为实。

要使 $\Psi$ 成为复数，必须离开 $|t|=1$。上半平面与下半平面中的 $\Psi$ 值都是可能的：前者对应于模大于 1 的实 $t$，后者对应于模小于 1 的实 $t$。

若 $t=\pm i$，则 $\Psi=\infty$，与 $\tau$ 无关。因此，在 $t=\pm i$ 处，$\tau$ 的取值完全无关紧要。这一点本质上已由 Marcus \cite{ref53} 直接论证。

最后一个观察：$\Psi$ 的一切值都可以在任意弱的耦合——即任意大的 ${\rm Im}\,\tau$——下实现。确实，把 $\Psi$ 的定义反解出来，我们有
% 式 (3.52)
\setcounter{equation}{51}
\begin{equation}\label{eq:helmo}
t^2=-{\Psi-\bar\tau\over \Psi-\tau},
\end{equation}
这表明，若固定 $\Psi$ 与 ${\rm Re}\,\tau$ 并令 ${\rm Im}\,\tau\to\infty$，则 $t\to\pm 1$。极限 ${\rm Im}\,\tau\to\infty$ 是弱耦合极限，此时理论应当是容易理解的。

\par\medskip\emph{$S$-对偶在典范参数上的作用}

现在我们要确定 $\Psi$ 在 $S$-对偶下如何变换。我们断言，使 $\tau$ 按 $\tau\to(a\tau+b)/(c\tau+d)$ 变换的变换以同样的方式作用于 $\Psi$：% 式 (3.53)
\setcounter{equation}{52}
\begin{equation}\label{eq:yrto}
\Psi\to
{a\Psi+b\over c\Psi+d}.
\end{equation}

显然，在保持 $t$ 不变的 $T:\tau\to\tau+1$ 下，我们有 $\Psi\to\Psi+1$。我们还必须检查在 $S:\tau\to-1/n_{\frak g}\tau$ 下的行为。按式~\eqref{eq:duco}，这一映射把 $t$ 变为 $\pm t\tau/|\tau|$。（符号取决于 $S$ 如何提升到 $SL(2,\Bbb{R})$；由于 $\Psi$ 在 $t\to-t$ 下不变，这无关紧要。）于是 $\Psi$ 按
% 式 (3.54)
\setcounter{equation}{53}
\begin{equation}\label{eq:nelmo}
\Psi\to -{1\over 2}\left({1\over n_{\frak
g}\tau}+{1\over n_{\frak g}\bar \tau}\right)-{1\over
2}\left({1\over n_{\frak g}\tau}-{1\over n_{\frak
g}\bar\tau}\right)\left({t^2(\tau/ \bar\tau)-1\over t^2(
\tau/\bar\tau)+1}\right)=-{1\over n_{\frak g} \Psi},
\end{equation}
变换，其中已用到式~\eqref{eq:helmo}。

\par\medskip{\it 定向反转}

发现像 $SL(2,\Bbb{Z})$ 这样的离散群作用于一个复参数 $\Psi$——如此处这样，$\Psi$ 取值于 $\Bbb{CP}^1$ 而不只是复平面的上半部分——有点不寻常。既然情况如此，就可以设想通过增加一个额外的对称性把 $SL(2,\Bbb{Z})$ 扩充为 $GL(2,\Bbb{Z})$，这个对称性可以取为% 式 (3.55)
\setcounter{equation}{54}
\begin{equation}\label{eq:ysnon}
{\EUT}=\left(\begin{matrix}-1 & 0 \\
0 & 1\end{matrix}\right).
\end{equation}
这样的元素将按 $\Psi\to-\Psi$ 作用于 $\Psi$，把上半平面与下半平面互换。

事实上，只要四维流形 $M$ 具有一个反转其定向的对称性 ${\EUT}:M\to M$，这样的对称性就存在。这样的变换保持规范耦合 $e^2$ 不变，从而保持 ${\rm Im}\,\tau=4\pi/e^2$ 不变；另一方面，它反转瞬子数的符号，从而反转 ${\rm Re}\,\tau$ 的符号。回顾式~\eqref{eq:yorgo} 中典范参数的定义可见，若 ${\EUT}$ 按% 式 (3.56)
\setcounter{equation}{55}
\begin{equation}\label{eq:zongor}
t\to
{1\over t}.
\end{equation}
作用，我们便得到 $\Psi\to-\Psi$。点 $t=\pm 1$ 是该变换的不动点。

\begin{yizhu}[含糊点拨]
按：定向反转把 $\Psi$ 换成 $-\Psi$，要同时看两件事。其一，$\Theta$ 保持 $e^2$（故 ${\rm Im}\,\tau$ 不变）但反转瞬子数的符号，等效于 $\tau\to-\bar\tau$，即 ${\rm Re}\,\tau$ 变号而 ${\rm Im}\,\tau$ 不变。其二，(3.56) 令 $t\to t^{-1}$，于是 (3.53) 中的 $(t-t^{-1})/(t+t^{-1})$ 变号。两者合起来，由本片译文所给的分解式 $\Psi={\rm Re}\,\tau+i\,{\rm Im}\,\tau\,(t-t^{-1})/(t+t^{-1})$ 立即得到 $\Psi\to-\Psi$。这正是 $\Psi\to-\Psi$ 可以作为理论的一个对称、从而把 $\Psi$ 的取值范围看成整个 $\Bbb{CP}^1$ 的原因。
\end{yizhu}

若取 ${\EUT}$ 按
% 式 (3.57)
\setcounter{equation}{56}
\begin{equation}\label{eq:yoro}
(A,\phi)\to ({\EUT}^*A,-{\EUT}^*\phi).
\end{equation}
作用于玻色子 $(A,\phi)$，我们便得到拓扑方程 \eqref{eq:loopme} 的一个对称性。对费米子的作用为% 式 (3.58)
\setcounter{equation}{57}
\begin{equation}\label{eq:zoro}
\begin{aligned}(\psi,\tilde\psi)&\to
({\EUT}^*\tilde\psi,{\EUT}^*\psi)\\ (\chi^+,\chi^-)&\to
({\EUT}^*\chi^-,{\EUT}^*\chi^+)\\ (\eta,\tilde\eta)&\to
({\EUT}^*\tilde\eta,{\EUT}^*\eta).\\\end{aligned}
\end{equation}
这给出作用量的一个显式对称性。超对称生成元按 $(u,v)\to(v,u)$ 变换，从而 $\delta_t\to t\delta_{t^{-1}}$。

一个重要的应用（我们将会看到）是 $M=\Bbb{R}\times W$ 的情形，其中 ${\EUT}$ 是「时间反演」对称性：它平凡地作用于 $W$，而在 $\Bbb{R}$ 上按乘以 $-1$ 的方式作用。
